%======================= MATERIAIS E MÉTODOS =======================
% Capítulo 4 — descreve a classificação metodológica, o hardware-alvo, a stack
% de ferramentas e, sobretudo, o PROTOCOLO DE AVALIAÇÃO (quatro métricas
% separadas + SUS + energia + BOM). Os resultados numéricos consolidados vão no
% Capítulo \ref{cap:resultados}; aqui citam-se apenas os valores de um piloto
% que JUSTIFICAM as escolhas de instrumentação.
\section{Materiais e métodos}\label{cap:metodos}

Este capítulo descreve como a prova de conceito (PoC) do DOMUS foi construída e,
principalmente, como ela é avaliada. A ênfase recai sobre o protocolo de avaliação:
para um trabalho cuja tese é que a democratização da automação residencial depende
mais da redução da barreira de comissionamento do que do preço do hardware, é o
método de medição --- e não a mera demonstração de funcionamento --- que sustenta ou
refuta a hipótese. A apresentação segue a NBR~14724 (ABNT, 2011) quanto à estrutura
do trabalho acadêmico.

\subsection{Classificação da pesquisa}\label{sec:classificacao}

Quanto à natureza, esta é uma \textbf{pesquisa aplicada}: parte de um problema concreto
--- famílias brasileiras que possuem hardware barato de automação, mas não conseguem
transformá-lo em automação utilizável --- e produz um artefato que endereça esse
problema. Quanto ao objetivo, caracteriza-se como \textbf{desenvolvimento de uma prova
de conceito (PoC)}: o produto é um sistema executável, o \textit{hub} DOMUS, cuja função
é evidenciar a viabilidade técnica do controle local por voz em português sem
dependência permanente de nuvem, e não um produto comercial acabado.

Metodologicamente, o desenvolvimento é \textbf{orientado a especificações}: cada
integração de dispositivo parte da leitura da documentação e do protocolo real do
fabricante (TP-LINK, 2024; TUYA, 2025) antes da escrita de código, evitando suposições
sobre APIs de IoT que mudam entre versões. Como estratégia central de projeto contra a
\textbf{fragmentação de protocolos} --- cada fabricante expõe um protocolo, um formato de
autenticação e um vocabulário de comandos distintos ---, adota-se o \textbf{padrão de
projeto \textit{Adapter}}. Nenhum controlador ou serviço de aplicação invoca as
bibliotecas de IoT diretamente; toda comunicação atravessa uma interface única
(\texttt{DeviceAdapter}), de modo que a heterogeneidade de Tuya e Tapo fique confinada a
implementações concretas e intercambiáveis. Essa decisão também viabiliza o \textbf{modo
MOCK}, no qual o sistema inteiro executa sem hardware físico, requisito para testes
automatizados reprodutíveis e para a integração contínua.

\subsection{Hardware-alvo}\label{sec:hardware}

O hardware-alvo foi selecionado para ser representativo do que uma família brasileira
compra em varejo popular: dispositivos Wi-Fi de 2,4~GHz, sem \textit{hub} proprietário
dedicado, de dois ecossistemas dominantes e de protocolos deliberadamente distintos, a
fim de exercitar o padrão \textit{Adapter} sob fragmentação real. A lâmpada Intelbras
pertence à linha Izy, que é \textit{white-label} do ecossistema Tuya e, portanto,
coberta pela mesma pilha de protocolo (TUYA, 2025); a tomada TP-Link Tapo utiliza o
protocolo proprietário KLAP (\textit{Key Local Authentication Protocol}). A
Tabela~\ref{tab:hardware} resume os dispositivos e a meta de custo total do arranjo,
fixada em menos de R\$~200 em hardware.

\begin{table}[H]
\centering
\caption{Hardware-alvo da prova de conceito e capacidades expostas.}
\label{tab:hardware}
\small
\begin{tabular}{@{}p{2.9cm}p{2.4cm}p{2.6cm}p{4.0cm}@{}}
\toprule
\textbf{Dispositivo} & \textbf{Modelo} & \textbf{Interface / protocolo} & \textbf{Capacidades expostas} \\
\midrule
Lâmpada LED & Intelbras EWS~410 & Wi-Fi 2,4~GHz / Tuya & liga/desliga, brilho, cor RGB \\
Tomada inteligente & TP-Link Tapo~P110 & Wi-Fi 2,4~GHz / KLAP & liga/desliga, medição de energia (W, kWh) \\
\midrule
\multicolumn{4}{@{}l}{\textit{Meta de custo do arranjo: total de hardware $<$ R\$~200,00 (verificado na BOM datada, Tab.~\ref{tab:bom-proto}).}} \\
\bottomrule
\end{tabular}
\end{table}

A escolha desses dois dispositivos não é arbitrária quanto à validação da tese. A tomada
Tapo~P110 permite exercitar o \textbf{caminho feliz} do controle local: seu protocolo
KLAP é acessível na LAN após uma autenticação local, e a medição de energia fornece um
sinal físico verificável de que um comando surtiu efeito. Já a lâmpada EWS~410 expõe a
\textbf{barreira de comissionamento} sob estudo: o pareamento inicial exige uma conta em
aplicativo proprietário e a obtenção de uma chave local (\texttt{local\_key}) via nuvem,
etapa que se mostrou, empiricamente, um ponto de bloqueio (Seção~\ref{sec:ambiente}).

\subsection{Ferramentas e tecnologias}\label{sec:ferramentas}

O \textit{hub} DOMUS é implementado em \textbf{Node.js~20~LTS} com \textbf{TypeScript} e
o \textit{framework} \textbf{NestJS}, cuja arquitetura modular por injeção de dependência
favorece o isolamento dos \textit{adapters}. A persistência usa \textbf{Prisma} sobre
\textbf{PostgreSQL}, e a comunicação em tempo real entre \textit{hub} e clientes usa
\textbf{Socket.IO} (WebSocket), de modo que mudanças de estado dos dispositivos sejam
refletidas na interface sem \textit{polling}. A interface de usuário é um
\textbf{\textit{Progressive Web App} (PWA)} construído em \textbf{Next.js~15} com
React~19, instalável no celular (\textit{Add to Home Screen}), o que dispensa publicação
em lojas de aplicativos e reforça a proposta \textit{local-first} (KLEPPMANN \textit{et
al.}, 2019).

O reconhecimento de fala (\textit{Speech-to-Text}, STT) usa o modelo \textbf{Whisper}
(RADFORD \textit{et al.}, 2022) executado \textbf{no próprio \textit{hub}} --- nunca no
celular e nunca na nuvem ---, decisão que é ao mesmo tempo de desempenho e de
privacidade. A integração com os dispositivos usa as bibliotecas \textbf{\texttt{tuyapi}}
(ecossistema Tuya/Intelbras) e \textbf{\texttt{tp-link-tapo-connect}} (Tapo), ambas
encapsuladas atrás da interface \texttt{DeviceAdapter}. A Tabela~\ref{tab:stack} consolida
as tecnologias por camada e sua função na PoC.

\begin{table}[H]
\centering
\caption{Ferramentas e tecnologias por camada.}
\label{tab:stack}
\small
\begin{tabular}{@{}p{3.0cm}p{5.0cm}p{4.2cm}@{}}
\toprule
\textbf{Camada} & \textbf{Tecnologia} & \textbf{Função na PoC} \\
\midrule
\textit{Hub}/backend & Node.js~20, TypeScript, NestJS & orquestração, regras, API \\
Persistência & Prisma + PostgreSQL & dispositivos, comandos, energia \\
Tempo real & Socket.IO (WebSocket) & estado ao vivo na interface \\
Interface & Next.js~15 + React~19 (PWA) & \textit{dashboard}, rotinas, voz \\
Voz (STT) & Whisper (no \textit{hub}) & transcrição de comandos pt-BR \\
IoT (\textit{adapters}) & \texttt{tuyapi}, \texttt{tp-link-tapo-connect} & controle local de Tuya e Tapo \\
\bottomrule
\end{tabular}
\end{table}

% Versões exatas conferem com apps/api/package.json: NestJS 11, Prisma 6,
% socket.io 4.8, tuyapi 7.7, tp-link-tapo-connect 2.0.15, nodejs-whisper 0.3.

\subsection{Protocolo de avaliação}\label{sec:protocolo}

A avaliação separa deliberadamente aquilo que a literatura de assistentes de voz
frequentemente confunde num único número. Uma única ``taxa de acerto'' agregada esconde
onde uma cadeia voz~$\rightarrow$~ação falha: se o erro está na transcrição, na
interpretação da intenção ou na execução física no dispositivo. Por isso, definem-se
\textbf{quatro métricas independentes} e três instrumentos complementares (SUS, energia e
BOM). A necessidade dessa separação é empiricamente motivada: em uma execução-piloto com
138 comandos reais, a acurácia de \emph{intenção} foi de 87,7\%, enquanto o sucesso de
\emph{execução} ponta-a-ponta foi de apenas 66,7\% --- diferença dominada por hardware
\textit{offline} no momento do comando, e não por falha do interpretador. Medir os dois
como um só número teria atribuído ao \textit{software} de linguagem uma falha que é do
dispositivo.

\subsubsection{Métrica 1 --- WER (taxa de erro de palavra)}

Avalia isoladamente a camada de transcrição do Whisper. Calcula-se a
\textit{Word Error Rate} pela razão entre edições e o total de palavras de referência:
\[
\mathrm{WER} = \frac{S + D + I}{N},
\]
em que $S$, $D$ e $I$ são, respectivamente, substituições, remoções e inserções em
relação à transcrição de referência, e $N$ é o número de palavras da referência. O
\textbf{corpus de avaliação} conterá \textbf{no mínimo 50 enunciados} de comando,
gravados por \textbf{ao menos 3 falantes} distintos, contemplando variação de sotaque e
formas coloquiais (por exemplo, ``desligua a luz'' além de ``desligar a luz''). Além do
WER agregado, produz-se uma \textbf{matriz de confusão} de comandos foneticamente
próximos, para identificar substituições sistemáticas (por exemplo, confusão entre
``ligar'' e ``desligar'' quando o áudio é curto). Registram-se também \textbf{alucinações}
do modelo --- transcrições plausíveis, mas sem fala correspondente ---, como o caso
observado em piloto em que ruído ambiente foi transcrito como ``[Som de futebol]''.
% A COLETAR: valor de WER agregado + matriz de confusão. Meta a definir após piloto.
O valor consolidado de WER é reportado no Capítulo~\ref{cap:resultados}; até o
fechamento deste texto permanece [a coletar].

\subsubsection{Métrica 2 --- acurácia de intenção}

Mede a camada de interpretação (\textit{parser}), isto é, se o texto transcrito foi
mapeado para a intenção correta (dispositivo + ação + parâmetros), \emph{independentemente}
de o dispositivo ter executado. Define-se como a razão entre comandos com intenção
corretamente classificada e o total de comandos, registrada na tabela
\texttt{voice\_commands} do banco. A \textbf{meta de projeto é $\geq 88\%$}; o piloto de
138 comandos já atingiu \textbf{87,7\%}, valor que baliza o parser e delimita o esforço
restante.

\subsubsection{Métrica 3 --- sucesso de execução ponta-a-ponta}

Mede a cadeia completa: um comando falado é bem-sucedido apenas se o dispositivo físico
atinge o estado pretendido, verificado por leitura de estado (para a tomada, confirmação
via medição de energia). É, portanto, a métrica mais severa e a que mais depende de
fatores fora do \textit{software} de linguagem. No piloto, o sucesso foi de
\textbf{66,7\%}, \textbf{dominado por dispositivos \textit{offline}} no instante do
comando (e não por erro de \textit{parsing}) --- distinção que só a separação em relação
à Métrica~2 torna visível. Falhas são categorizadas por causa (dispositivo \textit{offline},
\textit{timeout} de rede, erro de intenção) para diagnóstico.

\subsubsection{Métrica 4 --- latência voz$\rightarrow$execução (p50/p95)}

Mede o tempo entre o fim do enunciado e a confirmação de execução, reportado por
\textbf{percentis} (p50 e p95) em vez de média, por a distribuição ser assimétrica e de
cauda longa. O campo \texttt{latencyMs} é gravado por comando. No piloto, obteve-se
\textbf{p50~$=$~463~ms} e \textbf{p95~$=$~2140~ms}, com um \textit{outlier} máximo de
\textbf{12,3~s} atribuído a \textit{timeout} de hardware (dispositivo lento a responder),
e não a lentidão do \textit{hub}. Como referência de controle puramente local, medições
diretas da tomada Tapo~P110 via KLAP situaram-se entre \textbf{201 e 332~ms}, o que
estabelece o piso de latência da camada de dispositivo.

\begin{table}[H]
\centering
\caption{Resumo das quatro métricas do \textit{pipeline} de voz.}
\label{tab:metricas}
\small
\begin{tabular}{@{}p{0.5cm}p{3.0cm}p{4.4cm}p{3.4cm}@{}}
\toprule
\textbf{\#} & \textbf{Métrica} & \textbf{Definição / instrumento} & \textbf{Piloto / meta} \\
\midrule
1 & WER & $(S{+}D{+}I)/N$; corpus $\geq$50 enunciados, $\geq$3 falantes, matriz de confusão & [a coletar] \\
2 & Acurácia de intenção & intenções corretas $/$ total; tabela \texttt{voice\_commands} & 87,7\% (meta $\geq$88\%) \\
3 & Sucesso de execução & estado físico atingido $/$ total; leitura de estado & 66,7\% \\
4 & Latência voz$\to$exec. & percentis de \texttt{latencyMs} & p50~463~ms; p95~2140~ms \\
\bottomrule
\end{tabular}
\end{table}

\subsubsection{Usabilidade (SUS)}

A usabilidade percebida é medida pelo questionário \textit{System Usability Scale} (SUS),
de 10 itens em escala Likert de cinco pontos, aplicado a \textbf{$n = 3$ a $5$
participantes} após uma sessão guiada. Dada a natureza de PoC e o $n$ reduzido, os
resultados de SUS são tratados como \textbf{indicativos qualitativos}, não como inferência
estatística. Cada participante executa o \textbf{roteiro de 5 tarefas} da
Tabela~\ref{tab:roteiro-sus}, representativo do ciclo de uso doméstico, do comissionamento ao
acompanhamento de consumo. % Se o roteiro incorporar elementos de engajamento/onboarding, DETERDING et al. (2011) pode ser citado aqui.

\begin{table}[H]
\centering
\caption{Roteiro de 5 tarefas para a avaliação de usabilidade (SUS).}
\label{tab:roteiro-sus}
\small
\begin{tabular}{@{}p{0.6cm}p{12.4cm}@{}}
\toprule
\textbf{\#} & \textbf{Tarefa} \\
\midrule
T1 & Instalar o PWA no celular (\textit{Add to Home Screen}) e autenticar-se. \\
T2 & Adicionar/parear um dispositivo descoberto na rede local. \\
T3 & Ligar e desligar a lâmpada por comando de voz em português. \\
T4 & Criar uma automação por horário (ex.: desligar a tomada às 23h). \\
T5 & Consultar o consumo de energia da tomada no \textit{dashboard}. \\
\bottomrule
\end{tabular}
\end{table}

\subsubsection{Consumo de energia (antes/depois)}

Aproveitando a medição nativa da tomada Tapo~P110, registra-se o consumo (W instantâneo e
kWh acumulado) de uma carga controlada \textbf{antes} e \textbf{depois} da introdução de
uma automação (por exemplo, desligamento programado de um dispositivo em \textit{standby}),
na tabela \texttt{energy\_readings}. O objetivo é evidenciar o benefício prático da
automação, não uma medição metrológica de precisão. % A COLETAR: leituras W/kWh antes e depois; janela de observação de [...] dias.

\subsubsection{Custo (BOM datada)}

O custo total do hardware é documentado em uma \textbf{\textit{Bill of Materials} (BOM)
datada} --- Tabela~\ref{tab:bom-proto} ---, pois preços de varejo variam no tempo; a data de
coleta é parte do dado. A BOM verifica a meta de custo ($<$ R\$~200); estimativas
correntes situam o arranjo em torno de R\$~130, a confirmar na coleta final.

\begin{table}[H]
\centering
\caption{BOM datada do hardware-alvo (a coletar).}
\label{tab:bom-proto}
\small
\begin{tabular}{@{}p{4.2cm}p{1.4cm}p{2.4cm}p{2.4cm}p{2.0cm}@{}}
\toprule
\textbf{Item} & \textbf{Qtd.} & \textbf{Preço unit. (R\$)} & \textbf{Fonte/loja} & \textbf{Data} \\
\midrule
Intelbras EWS~410 & 1 & [...] & [...] & [dd/mm/aaaa] \\
TP-Link Tapo~P110 & 1 & [...] & [...] & [dd/mm/aaaa] \\
\midrule
\multicolumn{4}{@{}l}{\textbf{Total (meta $<$ R\$~200,00)}} & \textbf{[...]} \\
\bottomrule
\end{tabular}
\end{table}
% A COLETAR: preencher preços, lojas e datas reais de compra. Confirmar total < R$200.

\paragraph{Modo de avaliação e proteção de dados.}
Em operação normal, o DOMUS \textbf{não persiste áudio} de voz nem transcrições, em
coerência com o princípio de minimização de dados. Para computar WER e acurácia de
intenção, entretanto, é necessário reter transcrições comparáveis à referência. Isso é
feito por meio de um \textbf{``modo de avaliação''} explícito, ativado apenas durante as
sessões de coleta, no qual as \textbf{transcrições em texto} (não o áudio bruto) são
retidas \textbf{temporariamente} e apagadas após a análise. A ativação é precedida de
\textbf{consentimento informado} dos participantes, com finalidade, prazo de retenção e
procedimento de descarte declarados, em conformidade com a Lei Geral de Proteção de Dados
(BRASIL, 2018). Nenhum áudio é enviado à nuvem em nenhum modo de operação.

\subsection{Ambiente de testes}\label{sec:ambiente}

Toda a avaliação ocorre em uma \textbf{rede local (LAN)} doméstica de 2,4~GHz, sem
qualquer trânsito de áudio para a nuvem: o \textit{hub} DOMUS, os dispositivos e o celular
com o PWA residem no mesmo segmento de rede, e o reconhecimento de fala executa no
\textit{hub}. Essa topologia é a própria condição experimental da tese \textit{local-first}
(KLEPPMANN \textit{et al.}, 2019): mede-se o sistema sob a hipótese de operação sem nuvem.
% A COLETAR: especificar o hardware do hub (ex.: mini-PC/SBC, CPU, RAM) e o roteador.

Registra-se aqui uma \textbf{observação de ambiente empiricamente relevante} para os
resultados: o controle local da tomada Tapo~P110 via KLAP foi validado com sucesso
(latências de 201--332~ms), ao passo que a lâmpada Intelbras EWS~410 apresentou
\textbf{bloqueio no comissionamento} --- o pareamento inicial e a extração da
\texttt{local\_key} pela via de nuvem/aplicativo não se completaram no ambiente de teste.
Longe de ser um mero contratempo operacional, esse bloqueio constitui \textbf{evidência
empírica direta da barreira de comissionamento} que este trabalho investiga; sua
caracterização completa está \textbf{em validação} e é retomada no
Capítulo~\ref{cap:resultados}. A separação entre um caminho de controle que funciona
(Tapo) e uma barreira de entrada que trava (EWS~410) é, ela própria, um resultado do
método aqui descrito.
